% Content of the blueprint for the rumor-spreading (uniform push) formalization.
% Mirrors latex/rumor_push.tex lemma-for-lemma, but every statement here also
% carries \lean{...} / \leanok / \uses{...} tags so that leanblueprint can
% build the dependency graph and check coverage against RumorSpread/*.lean.

\begin{abstract}
\noindent
We give a complete and elementary proof that the uniform \emph{push} protocol
on the complete graph $K_n$, started from a single informed node, informs all
$n$ nodes within $O(\log n)$ rounds with high probability. The probability
space is finite and uniform throughout; the only analytic ingredients are
Bernoulli's inequality and Markov's inequality, and the only concentration
tool is a one-line exponential-moment induction for adaptive Bernoulli
trials. No measure theory, no PMF/\texttt{ENNReal}, no Chernoff bounds, and
no martingales are used. This argument has been \emph{fully formalized} (no
\texttt{sorry}) in Lean~4 + Mathlib; the main theorem is
\texttt{RumorPush.push\_informs\_all\_whp}.
\end{abstract}

\section{The model}\label{sec:model}

Fix $n \ge 2$ and let $V = \{0,1,\dots,n-1\}$ be the vertex set of the
complete graph $K_n$.

\begin{definition}[Round randomness]\label{def:round}
  \lean{RumorPush.Tgt}
  \leanok
  A \emph{round configuration} is a function
  $\omega \colon V \to V$ with $\omega(v) \ne v$ for all $v$.
  Let $R_n$ denote the (finite) set of round configurations; it carries the
  uniform probability measure, i.e.\ every node picks a uniformly random
  \emph{target} among the other $n-1$ nodes, independently of all other
  nodes.
\end{definition}

\begin{definition}[Push step]\label{def:step}
  \lean{RumorPush.step}
  \leanok
  \uses{def:round}
  For an \emph{informed set} $I \subseteq V$ and a round configuration
  $\omega \in R_n$, the next informed set is
  \[
    \mathrm{step}(I,\omega) \;=\; I \,\cup\, \omega(I)
    \;=\; I \cup \{\omega(v) : v \in I\}.
  \]
  (Only informed nodes actually transmit; letting \emph{all} nodes draw a
  target and ignoring the choices of uninformed nodes is an equivalent and
  formalization-friendlier convention.)
\end{definition}

\begin{definition}[Process]\label{def:process}
  \lean{RumorPush.run}
  \leanok
  \uses{def:step}
  Fix a horizon $T \in \N$ and a start vertex $v_0 \in V$. The sample space
  is $\Omega = R_n^{\,T}$ (uniform measure, i.e.\ i.i.d.\ uniform rounds).
  For $\omega = (\omega_0,\dots,\omega_{T-1}) \in \Omega$ define
  $I_0(\omega) = \{v_0\}$ and
  $I_{t+1}(\omega) = \mathrm{step}(I_t(\omega), \omega_t)$ for $0 \le t < T$.
  We write $m_t = |I_t|$ and $U_t = n - m_t$ for the numbers of informed and
  uninformed nodes. Since $\Omega$ is finite, all probabilities are finite
  ratios and all expectations are finite sums; ``conditioning on $I_t$''
  means partitioning $\Omega$ according to the first $t$ coordinates. In the
  formalization, a trajectory is simply a \texttt{List} of round
  configurations and $\mathrm{run}\ I\ l$ folds \texttt{step} over the list
  $l$.
\end{definition}

Two trivial but load-bearing facts, both immediate from
Definition~\ref{def:step}:

\begin{lemma}[Monotonicity]\label{lem:mono}
  \lean{RumorPush.subset_step, RumorPush.card_le_card_step,
    RumorPush.subset_run}
  \leanok
  \uses{def:process}
  $I_t(\omega) \subseteq I_{t+1}(\omega)$ for all $t$ and all $\omega$;
  hence $m_t$ is non-decreasing and $U_t$ is non-increasing in $t$,
  pointwise on $\Omega$.
\end{lemma}

\begin{lemma}[At most doubling]\label{lem:atmost}
  \lean{RumorPush.card_step_le}
  \leanok
  \uses{def:step}
  $|I_{t+1}(\omega) \smallsetminus I_t(\omega)| \le |I_t(\omega)|$; equivalently
  $(\mathrm{step}(I,\omega)).\mathrm{card} \le 2\,I.\mathrm{card}$.
\end{lemma}

\section{Main theorem}\label{sec:main-statement}

Throughout, $\log$ denotes the natural logarithm.

\begin{theorem}[Push protocol informs $K_n$ in $O(\log n)$ rounds w.h.p.]
  \label{thm:main}
  \lean{RumorPush.push_informs_all_whp, RumorPush.push_informs_all_whp'}
  \leanok
  \uses{lem:gluing, lem:constants, lem:phase2cor}
  Let $n \ge 2$, let $v_0 \in V$, and set
  \[
    T_1 \;=\; \ceil{117 \log n} + 23, \qquad
    T_2 \;=\; \ceil{6 \log n}, \qquad
    T \;=\; T_1 + T_2 .
  \]
  Then $\Pr\bigl[\, I_T = V \,\bigr] \;\ge\; 1 - \frac{2}{n}$.
  In particular all nodes are informed after $O(\log n)$ rounds with high
  probability.
\end{theorem}

\begin{remark}
  The constants are deliberately crude: the proof is organized to minimize
  the analytic toolkit, not the round count. The sharp answer is
  $\log_2 n + \log n + o(\log n)$ (Frieze--Grimmett \cite{FG85}, Pittel
  \cite{Pittel87}). Replacing $2/n$ by $n^{-\alpha}$ for any fixed
  $\alpha \ge 1$ only changes the constants in $T_1, T_2$.
\end{remark}

The proof has two phases, glued by Lemma~\ref{lem:mono}:
\begin{itemize}
\item \textbf{Growth phase} (Section~\ref{sec:growth}): while $m_t \le n/2$,
  each round multiplies $m_t$ by $\ge 9/8$ with probability $\ge 1/8$;
  an exponential-moment bound for adaptive trials shows that after $T_1$
  rounds, $m_{T_1} > n/2$ except with probability $\le 1/n$.
\item \textbf{Saturation phase} (Section~\ref{sec:saturation}): once
  $m_t \ge n/2$, the expected number of uninformed nodes contracts by a
  factor $2/3$ per round; after $T_2$ further rounds Markov's inequality
  gives $U_T = 0$ except with probability $\le 1/n$.
\end{itemize}

\section{Finite uniform probability}\label{sec:prelim}

All randomness in this development is uniform over finite types, so instead
of measure theory or \texttt{PMF} we work with plain finite sums.

\begin{definition}[Average]\label{def:avg}
  \lean{RumorPush.avg}
  \leanok
  For a finite type $\alpha$ and $f \colon \alpha \to \R$,
  $\mathrm{avg}(f) := \bigl(\sum_{a} f(a)\bigr) / |\alpha|$ is the
  expectation of $f$ under the uniform distribution on $\alpha$.
\end{definition}

\begin{definition}[Trajectory expectation]\label{def:expList}
  \lean{RumorPush.expList}
  \leanok
  \uses{def:avg}
  For a finite type $\alpha$, $\mathrm{expList}\ \alpha\ T\ F$ is the
  expectation of a functional $F$ of a trajectory of $T$ i.i.d.\ uniform
  draws from $\alpha$, defined \emph{by recursion on $T$}: the head draw is
  averaged out first, i.e.\ $\mathrm{expList}\ \alpha\ 0\ F = F(\,[]\,)$ and
  $\mathrm{expList}\ \alpha\ (T{+}1)\ F
    = \mathrm{avg}\bigl(a \mapsto \mathrm{expList}\ \alpha\ T\ (l \mapsto
    F(a :: l))\bigr)$.
  This makes ``conditioning on the first round'' a definitional unfolding
  rather than a measure-theoretic construction.
\end{definition}

\begin{lemma}[Basic properties]\label{lem:avg-basic}
  \lean{RumorPush.avg_le_avg, RumorPush.avg_add, RumorPush.avg_sub,
    RumorPush.avg_const_mul, RumorPush.avg_const,
    RumorPush.expList_le_expList, RumorPush.expList_add,
    RumorPush.expList_const_mul, RumorPush.expList_const}
  \leanok
  \uses{def:avg, def:expList}
  $\mathrm{avg}$ and $\mathrm{expList}$ are monotone, additive, and
  homogeneous, and constants pass through unchanged. These pointwise
  inequalities pushed through $\mathrm{avg\_le\_avg}$ /
  $\mathrm{expList\_le\_expList}$ are the only form in which Markov's
  inequality and union bounds are used below (Lemmas~\ref{lem:markov}
  and~\ref{lem:revmarkov}); they are not separately formalized lemmas.
\end{lemma}

\begin{theorem}[Faithfulness of the model]\label{thm:faithful}
  \lean{RumorPush.expList_eq_avg_ofFn}
  \leanok
  \uses{def:expList}
  $\mathrm{expList}\ \alpha\ T\ F$ equals the uniform average of $F$ over
  \emph{all} length-$T$ sequences of draws, i.e.\ over the product
  probability space $\mathrm{Fin}\ T \to \alpha$ of $T$ i.i.d.\ uniform
  rounds. This confirms that the recursive definition of
  Definition~\ref{def:expList} is the textbook object.
\end{theorem}

\section{Elementary inequalities}\label{sec:ineq}

We write $x := \tfrac{1}{n-1}$, so $0 < x \le 1$ for $n \ge 2$.

\begin{lemma}[Bernoulli]\label{lem:bernoulli}
  \lean{RumorPush.one_sub_mul_le_pow}
  \leanok
  For all real $y \ge -1$ (we only use $y\in[-1,1]$) and $m \in \N$:
  $(1+y)^m \ge 1 + my$. Equivalently, with $y = -x$:
  $1 - mx \le (1-x)^m$.
\end{lemma}

\begin{lemma}[Second-order lower bound, ``Bonferroni'']\label{lem:bonferroni}
  \lean{RumorPush.mul_sub_sq_le_one_sub_pow, RumorPush.half_mul_le_one_sub_pow}
  \leanok
  \uses{lem:bernoulli}
  For $0 \le x \le 1$ and $m \in \N$:
  \[
    1-(1-x)^m \;\ge\; mx - \binom{m}{2}x^2 \;\ge\; mx\Bigl(1-\frac{mx}{2}\Bigr),
  \]
  and consequently $mx/2 \le 1-(1-x)^m$ whenever $mx \le 1$.
\end{lemma}

\begin{lemma}[Upper bound without $e$]\label{lem:upper}
  \lean{RumorPush.pow_one_sub_le_one_div}
  \leanok
  \uses{lem:bernoulli}
  For $0 \le x \le 1$ and $m \in \N$: $(1-x)^m \le \dfrac{1}{1+mx}$.
\end{lemma}

\begin{lemma}[Logarithm bounds]\label{lem:logbounds}
  \lean{RumorPush.one_sub_inv_le_log}
  \leanok
  For $x > 0$: $1 - \tfrac1x \le \log x$. (The formalization only needs this
  left inequality, applied at $x = 9/8,\ 16/15,\ 3/2$, together with
  Mathlib's numeric bound $\log 2 < 0.6931471808$.)
\end{lemma}

\begin{lemma}[Markov's inequality, finite uniform version]\label{lem:markov}
  \lean{RumorPush.avg_le_avg}
  \leanok
  \uses{def:avg}
  Let $\Omega$ be a finite nonempty set with the uniform measure,
  $f \colon \Omega \to \R$ with $f \ge 0$, and $a > 0$. Then
  $\Pr[f \ge a] \le \E[f]/a$. In particular, if $f$ takes values in $\N$,
  then $\Pr[f \ne 0] \le \E[f]$.
\end{lemma}

\begin{lemma}[Reverse Markov for bounded variables]\label{lem:revmarkov}
  \lean{RumorPush.avg_le_avg}
  \leanok
  \uses{def:avg}
  Let $f \colon \Omega \to \R$ with $0 \le f \le b$ pointwise, $b > 0$, and
  let $0 < a < b$. Then $\Pr[f \ge a] \ge \dfrac{\E[f] - a}{b}$. (As with
  Lemma~\ref{lem:markov}, the formalization does not isolate this as a
  named lemma: each use site supplies its own pointwise bound and pushes it
  through $\mathrm{avg\_le\_avg}$, e.g.\ in the proof of
  Lemma~\ref{lem:goodround}.)
\end{lemma}

\section{One-round estimates}\label{sec:oneround}

Fix an informed set $I$ with $|I| = m \ge 1$ and let $\omega \in R_n$ be a
uniform round configuration. All statements in this section concern this
single round; by independence of the rounds they apply verbatim to the
conditional law of round $t$ given $I_t = I$.

\begin{lemma}[Contact probability]\label{lem:contact}
  \lean{RumorPush.avg_not_contacted}
  \leanok
  \uses{def:step, def:avg}
  For any fixed $u \notin I$,
  \[
    \Pr_\omega\bigl[u \in \mathrm{step}(I,\omega)\bigr]
    \;=\; 1 - \Bigl(1-\frac{1}{n-1}\Bigr)^{m}.
  \]
  The Lean proof is the only place where a probability is \emph{computed}
  rather than estimated: it counts round configurations avoiding $u$ on $I$
  via an explicit equivalence to a product of subtypes and
  \texttt{Fintype.card\_pi}.
\end{lemma}

\begin{lemma}[Expected growth]\label{lem:expgrowth}
  \lean{RumorPush.avg_card_step}
  \leanok
  \uses{lem:contact, lem:bonferroni}
  Let $N(\omega) = |\mathrm{step}(I,\omega)| - m$ be the number of newly
  informed nodes. The formalization gives the \emph{exact} formula
  \[
    \E_\omega[|\mathrm{step}(I,\omega)|] \;=\;
      m + (n-m)\bigl(1-(1-x)^m\bigr),
  \]
  from which, if $1 \le m \le n/2$: $\E_\omega[N] \ge m/4$ (using
  Lemma~\ref{lem:bonferroni} and $n - m \ge n/2$, $mx \le 1$).
\end{lemma}

\begin{lemma}[A good round has constant probability]\label{lem:goodround}
  \lean{RumorPush.prob_goodRound}
  \leanok
  \uses{lem:expgrowth, lem:atmost, lem:revmarkov}
  If $1 \le m \le n/2$, then
  $\Pr_\omega\bigl[\,|\mathrm{step}(I,\omega)| \ge \tfrac{9}{8}\,m\,\bigr]
    \ge \tfrac18$.
\end{lemma}

\begin{lemma}[Contraction above half]\label{lem:contract}
  \lean{RumorPush.avg_uninformed_le}
  \leanok
  \uses{lem:contact, lem:upper}
  If $m \ge n/2$, then for every fixed $u \notin I$,
  $\Pr_\omega[u \notin \mathrm{step}(I,\omega)] = (1-x)^m \le \tfrac23$, and
  consequently, with $U' := n - |\mathrm{step}(I,\omega)|$,
  $\E_\omega[U'] \le \tfrac23(n-m)$.
\end{lemma}

\section{Adaptive Bernoulli trials}\label{sec:adaptive}

\begin{definition}[Good rounds]\label{def:good}
  \lean{RumorPush.goodRound, RumorPush.goodCount}
  \leanok
  \uses{def:step}
  A round is \emph{good} if it multiplies the informed set by $\ge 9/8$, or
  if the informed set already exceeds $n/2$:
  \[
    X_t \;=\; \mathbf{1}\Bigl[\, m_{t+1} \ge \tfrac98\, m_t
    \;\;\text{or}\;\; m_t > \tfrac n2 \,\Bigr] \in \{0,1\}.
  \]
  $\mathrm{goodCount}$ counts the good rounds along a trajectory. By
  Lemma~\ref{lem:goodround} (applied when $m_t \le n/2$; trivial otherwise),
  $\Pr[X_t = 1 \mid \omega_0,\dots,\omega_{t-1}] \ge \tfrac18$ for every
  prefix, since $\omega_t$ is independent of the prefix and
  $\mathrm{expList}$ conditions on it definitionally.
\end{definition}

\begin{lemma}[Exponential moment for adaptive trials]\label{lem:adaptive}
  \lean{RumorPush.expList_half_pow_goodCount}
  \leanok
  \uses{def:expList, def:good, lem:goodround}
  For every $T$ and every nonempty $I$,
  \[
    \E\Bigl[\bigl(\tfrac12\bigr)^{\mathrm{goodCount}(I,\,\cdot\,)}\Bigr]
      \;\le\; \Bigl(\frac{15}{16}\Bigr)^{T}
  \]
  over $T$ rounds started at $I$. (This specializes the textbook
  ``$\E[s^G] \le (1-p(1-s))^T$ for adaptive $\{0,1\}$-trials with
  per-step success probability $\ge p$'' at $p = 1/8$, $s = 1/2$; the
  formalization proves the specialized statement directly by induction on
  $T$ rather than the general one, since only this instance is needed. The
  ensuing tail bound $\Pr[G < L] \le 2^{L}(15/16)^{T}$ — via
  $s^{L}\Pr[G<L] \le \E[s^G]$ — is likewise not isolated as a separate
  lemma; it is proved inline inside Lemma~\ref{lem:phase1}
  (\texttt{phase1}).
\end{lemma}

\section{Growth phase}\label{sec:growth}

\begin{lemma}[Enough good rounds force half-saturation]\label{lem:force}
  \lean{RumorPush.pow_goodCount_mul_card_le_card_run}
  \leanok
  \uses{def:good, lem:mono}
  Let $L \in \N$ satisfy $(9/8)^L > n/2$. If the informed set never exceeds
  $n/2$ up to time $T_1$, then along that trajectory
  $m_{T_1} \ge (9/8)^{\mathrm{goodCount}} \cdot m_0$: every good round
  multiplies $m_t$ by $\ge 9/8$ and every round is non-decreasing
  (Lemma~\ref{lem:mono}), so $\mathrm{goodCount} \ge L$ would force
  $m_{T_1} > n/2$, a contradiction — hence
  $\mathrm{goodCount} < L \Rightarrow m_{T_1} \le n/2$ is impossible to rule
  out only via the deterministic growth bound stated here.
\end{lemma}

\begin{lemma}[Phase 1]\label{lem:phase1}
  \lean{RumorPush.phase1}
  \leanok
  \uses{lem:force, lem:adaptive}
  If $(9/8)^L \ge n$, then after $T_1$ rounds the informed set is still
  $\le n/2$ with probability at most $2^{L}\,(15/16)^{T_1}$.
\end{lemma}

\begin{lemma}[Admissible constants]\label{lem:constants}
  \lean{RumorPush.numeric_A, RumorPush.numeric_B}
  \leanok
  \uses{lem:phase1, lem:logbounds}
  With $L = \ceil{9 \log n} + 1$ we have $(9/8)^L \ge n$
  (\texttt{numeric\_A}), and with $T_1 = \ceil{117 \log n} + 23$ we have
  $2^{L}(15/16)^{T_1} \le 1/n$ (\texttt{numeric\_B}), using
  $\log\tfrac98 \ge \tfrac19$, $\log\tfrac{16}{15} \ge \tfrac1{16}$ and
  $\log 2 < 0.6932$. Combined with Lemma~\ref{lem:phase1}:
  $\Pr[m_{T_1} \le n/2] \le 1/n$. The clean statement to formalize is that
  \emph{any} $L, T_1$ satisfying these two numeric inequalities work; the
  values above merely exhibit admissible ones.
\end{lemma}

\section{Saturation phase}\label{sec:saturation}

\begin{lemma}[Phase 2]\label{lem:phase2}
  \lean{RumorPush.saturation}
  \leanok
  \uses{lem:contract}
  Let $A = \{m_{T_1} > n/2\}$ (an event determined by the first $T_1$
  coordinates). Then for every $s \ge 0$,
  $\E\bigl[\, U_{T_1 + s}\, \mathbf{1}_A \,\bigr] \le (2/3)^{s}\, n$, by
  induction on $s$ using Lemma~\ref{lem:contract} at each step (monotonicity
  keeps $m_{T_1+s} > n/2$ on $A$).
\end{lemma}

\begin{corollary}\label{lem:phase2cor}
  \lean{RumorPush.numeric_C}
  \leanok
  \uses{lem:phase2, lem:markov, lem:logbounds}
  With $T_2 = \ceil{6 \log n}$: $\Pr[\, U_T \ge 1 \text{ and } A \,]
    \le n\,(2/3)^{T_2} \le 1/n$. ($U_T$ is $\N$-valued, so Markov applied to
  $U_T \mathbf{1}_A$ and Lemma~\ref{lem:phase2} give the first inequality;
  $(2/3)^{T_2} n \le 1/n$ — \texttt{numeric\_C} — follows from
  $\log\tfrac32 \ge \tfrac13$ and $T_2 \ge 6\log n$.)
\end{corollary}

\section{Proof of the main theorem}\label{sec:proof-main}

\begin{lemma}[Gluing the two phases]\label{lem:gluing}
  \lean{RumorPush.prNotAllInformed, RumorPush.prNotAllInformed_le}
  \leanok
  \uses{lem:phase1, lem:phase2, lem:mono}
  Let $\mathrm{prNotAllInformed}(n,v_0,T)$ be the probability that some node
  is still uninformed after $T$ rounds started at $v_0$. For all
  $L, T_1, T_2$ with $(9/8)^L \le n$:
  \[
    \mathrm{prNotAllInformed}(n,v_0,T_1+T_2) \;\le\;
      2^{L}\Bigl(\frac{15}{16}\Bigr)^{T_1} + \Bigl(\frac23\Bigr)^{T_2} n.
  \]
  Failure requires either phase 1 to fail (informed set still $\le n/2$
  after $T_1$ rounds, Lemma~\ref{lem:phase1}) or phase 2 to leave an
  uninformed node on the complementary event, controlled by
  Lemma~\ref{lem:phase2} via Markov.
\end{lemma}

With $T_1 = \ceil{117\log n} + 23$, $T_2 = \ceil{6\log n}$, $T = T_1+T_2$,
Theorem~\ref{thm:main} follows from Lemma~\ref{lem:gluing},
Lemma~\ref{lem:constants} and Corollary~\ref{lem:phase2cor}:
\[
  \Pr[I_T \ne V] \;=\; \Pr[U_T \ge 1]
  \;\le\; \Pr[\,\overline{A}\,] + \Pr[\,U_T \ge 1 \text{ and } A\,]
  \;\le\; \frac1n + \frac1n \;=\; \frac2n .
\]

\section{The Lean formalization}\label{sec:lean}

The argument above is fully formalized in the accompanying repository
(library \texttt{RumorSpread}, namespace \texttt{RumorPush}); the build
contains no \texttt{sorry} and the main theorem depends only on the three
standard axioms (\texttt{propext}, \texttt{Classical.choice},
\texttt{Quot.sound}). File structure (dependency order):
\texttt{Bounds.lean} (Section~\ref{sec:ineq}), \texttt{Prob.lean}
(Section~\ref{sec:prelim}), \texttt{Model.lean}
(Section~\ref{sec:model}), \texttt{OneRound.lean}
(Section~\ref{sec:oneround}), \texttt{Growth.lean}
(Sections~\ref{sec:adaptive}--\ref{sec:growth}), \texttt{Saturation.lean}
(Section~\ref{sec:saturation}), \texttt{Main.lean} (main theorem and
numeric lemmas), \texttt{Equivalence.lean} (Theorem~\ref{thm:faithful}).

\begin{thebibliography}{9}
\bibitem{FG85} A.~Frieze, G.~Grimmett.
\emph{The shortest-path problem for graphs with random arc-lengths}.
Discrete Applied Mathematics 10 (1985) 57--77.
\bibitem{Pittel87} B.~Pittel.
\emph{On spreading a rumor}. SIAM J.\ Appl.\ Math.\ 47 (1987) 213--223.
\bibitem{KSSV00} R.~Karp, C.~Schindelhauer, S.~Shenker, B.~V\"ocking.
\emph{Randomized rumor spreading}. FOCS 2000.
\end{thebibliography}
